% Retrievalprestatie kabinetsreactie-matcher (DV1).
% Tekst geverifieerd tegen
% kabinetsreactie_data/retrieval_report.json (advies_count 69, precision@1 0,5652,
% recall@50 0,8841, MRR 0,6749).
% Aanvullende VLAM-poort-data (precision 0,64, VLAM-recall 0,97, end-to-end 0,81)
% staat in kabinetsreactie_data/vlam_precision_report.json; nog niet in deze tekst verwerkt.

\subsubsection{Retrievalprestatie van de kabinetsreactie-matcher.}
De kabinetsreactie-matcher werkt in twee fasen: een retrievallaag zoekt
kandidaat-documenten die mogelijk als reactie op een adviesrapport dienen,
waarna de VLAM-poort die kandidaten inhoudelijk beoordeelt
(Tabel~\ref{tab:dv1-matchers}). De hier gerapporteerde evaluatie betreft
uitsluitend de retrievallaag: de vraag of de juiste reactie \"uberhaupt in de
kandidaatlijst verschijnt, nog v\'o\'or de validatie die tot de 232 gevalideerde
koppelingen leidt. Alleen de deterministische zoekroutes stonden aan; de
semantische pgvector-route was uitgeschakeld, zodat de gerapporteerde recall een
ondergrens vormt.

\noindent Als grondwaarheid dienden advies-reactieparen waarin het adviescollege
zelf het betreffende Kamerstuk als kabinetsreactie publiceerde en waarin dat
doeldocument aantoonbaar in de doorzochte corpus aanwezig was. Dit is een aparte,
geconstrueerde benchmarkset en geen subset van de 417 geanalyseerde
adviesrapporten: het criterium identificeerde 183 paren over 167 adviesrapporten,
waarvan deze pilot er 69 evalueerde. Net als bij de instellingsbesluit-matcher
(candidate recall 119/151) gaat het dus om recall t\'egen bekende paren, niet om
volledige recall over alle werkelijke koppelingen
(zie Tabel~\ref{tab:dv1-beperkingen}).

\noindent De juiste reactie staat in 57\,\% van de gevallen direct op de eerste
positie, in 75\,\% binnen de top drie en in 88\,\% binnen de top vijftig. De mean
reciprocal rank (MRR), het gemiddelde van de reciproque rangposities, bedraagt
0,67; dat komt overeen met een juist document dat doorgaans hoog in de lijst staat,
rond positie 1 \`a 2 (Tabel~\ref{tab:dv1-retrieval}). Geen van de 69 gevallen
leverde een lege resultatenlijst op. Deze cijfers suggereren dat de retrievallaag
een solide basis vormt, maar dat de daaropvolgende VLAM-poort nodig blijft om
valse positieven te filteren en de kandidaten te prioriteren, passend bij de
bewust lage precisie van de brede zoekstap. De schattingen zijn omgeven met
onzekerheid vanwege de kleine, exploratieve steekproef ($N = 69$), en de
grondwaarheid uit door colleges zelf gepubliceerde reacties kan selectiebias
bevatten, zodat de cijfers niet zonder voorbehoud op de volledige populatie van
toepassing zijn.

\begin{table}[htbp]
\centering
\caption{Retrievalprestatie van de kabinetsreactie-matcher op een geconstrueerde
benchmarkset ($N = 69$ adviesrapporten), met 95\%-betrouwbaarheidsintervallen
volgens Wilson. Alleen deterministische zoekroutes; de semantische route stond
uit, zodat recall een ondergrens is.}
\label{tab:dv1-retrieval}
\footnotesize
\begin{tabular}{@{}lcc@{}}
\toprule
Metriek & Waarde & 95\,\%-BI \\
\midrule
Precision@1 & 0,57 & 0,45--0,68 \\
Recall@1    & 0,57 & 0,45--0,68 \\
F1 (top-1)  & 0,57 & --- \\
Recall@3    & 0,75 & 0,64--0,84 \\
Recall@10   & 0,87 & 0,77--0,93 \\
Recall@25   & 0,88 & 0,79--0,94 \\
Recall@50   & 0,88 & 0,79--0,94 \\
MRR         & 0,67 & --- \\
\bottomrule
\end{tabular}
\end{table}
